% ============================================================================
% TECHNICAL APPENDIX 7A: Annotation Metrics, Agreement Statistics, and Label Quality
% Formal treatment of inter-annotator agreement and annotation pipeline design
% ============================================================================
% Source: /markdown/ch7/Ch7A_final.md
% ============================================================================

\chapter*{Technical Appendix 7A: Annotation Metrics and Label Quality}
\addcontentsline{toc}{chapter}{Technical Appendix 7A: Annotation Metrics}

\begin{chapterquote}{Formal foundations for the examples in Chapter 7}
The annotation challenges described in Chapter~7 can be formalized using agreement statistics, Bayesian aggregation models, and signal detection theory.
\end{chapterquote}

% ============================================================================
\section{Cohen's Kappa: Derivation and Properties}
% ============================================================================

\subsection{Basic Formula}

For two annotators labeling $N$ items into $k$ categories, let $n_{ij}$ be the number of items where annotator~1 assigned category~$i$ and annotator~2 assigned category~$j$.

\[
P_o = \frac{\sum_i n_{ii}}{N} \quad \text{(observed agreement)}
\]

\[
P_e = \sum_i \frac{n_{i\cdot}}{N} \cdot \frac{n_{\cdot i}}{N} \quad \text{(expected chance agreement)}
\]

\[
\kappa = \frac{P_o - P_e}{1 - P_e}
\]

where $n_{i\cdot} = \sum_j n_{ij}$ is the row marginal and $n_{\cdot i} = \sum_j n_{ji}$ is the column marginal.

\subsection{Variance and Confidence Intervals}

\[
\text{Var}(\kappa) \approx \frac{P_o(1-P_o)}{N(1-P_e)^2}
\]

Approximate 95\% confidence interval: $\kappa \pm 1.96\sqrt{\text{Var}(\kappa)}$

\begin{cautionbox}
This variance formula is a simplified approximation that ignores variance from estimating $P_e$ itself, producing confidence intervals that are slightly too narrow---especially when $P_e$ is large or $N$ is small. For publication-quality intervals use the exact formula from \textcite{fleiss1971measuring} or a bootstrap estimate.
\end{cautionbox}

For reliable kappa estimates, require $N \geq 30$ items per annotator pair.

\subsection{Weighted Kappa for Ordinal Categories}

With linear weights $w_{ij} = 1 - |i-j|/(k-1)$ or quadratic weights $w_{ij} = 1 - [(i-j)/(k-1)]^2$, define weighted observed and expected agreement:

\[
P_o^w = \sum_{ij} \frac{w_{ij}\, n_{ij}}{N}, \qquad P_e^w = \sum_{ij} w_{ij}\, P_{e,ij},
\quad P_{e,ij} = \frac{n_{i\cdot}}{N} \cdot \frac{n_{\cdot j}}{N}
\]

Weighted kappa follows the same structure as unweighted kappa \autocite{cohen1968weighted}:

\[
\kappa_w = \frac{P_o^w - P_e^w}{1 - P_e^w}
\]

Quadratic weights are standard for medical rating scales; linear weights for quality assessments.

% ============================================================================
\section{Fleiss' Kappa: Extension to Multiple Annotators}
% ============================================================================

For $M$ annotators labeling $N$ items into $k$ categories, let $n_{ij}$ be the number of annotators who assigned item~$i$ to category~$j$:

\[
\bar{P} = \frac{1}{NM(M-1)} \sum_i \sum_j n_{ij}(n_{ij}-1)
\qquad
\bar{P}_e = \sum_j \left(\frac{\sum_i n_{ij}}{NM}\right)^2
\]

\[
\kappa_F = \frac{\bar{P} - \bar{P}_e}{1 - \bar{P}_e}
\]

\begin{table}[htbp]
\caption{Agreement benchmarks by annotation domain}
\label{tab:kappa-benchmarks}
\centering
\begin{tabular}{@{}lcc@{}}
\toprule
\textbf{Domain} & \textbf{Typical Range} & \textbf{Target} \\
\midrule
Binary toxicity (general)      & 0.45--0.65 & $\geq 0.60$ \\
Medical image (radiology)      & 0.55--0.75 & $\geq 0.75$ \\
Sentiment analysis (clear)     & 0.70--0.85 & $\geq 0.70$ \\
Named entity recognition       & 0.80--0.95 & $\geq 0.85$ \\
Legal document classification  & 0.50--0.70 & $\geq 0.65$ \\
RLHF instruction quality       & 0.55--0.75 & $\geq 0.65$ \\
AAVE toxicity (diverse pool)   & 0.40--0.55 & $\geq 0.60$ \\
\bottomrule
\end{tabular}
\end{table}

% ============================================================================
\section{Annotation as a Bayesian Inference Problem: Dawid-Skene}
% ============================================================================

\subsection{Generative Model}

The Dawid-Skene model \autocite{dawid1979maximum} treats annotator labels as noisy observations of a true latent label, estimating both true labels and annotator error rates jointly.

For item $i$, annotator $j$, and observed label $y_{ij}$:
\[
P(y_{ij} = l \mid z_i = k) = \pi_{jkl}
\]
where $z_i$ is the true label for item $i$ and $\pi_{jkl}$ is annotator $j$'s probability of labeling as $l$ when the true label is $k$.

\subsection{EM Algorithm}

\paragraph{E-step.} Compute posterior probability of true label:
\[
P(z_i = k \mid \mathbf{y}_i) \propto p_k \prod_j \prod_l \pi_{jkl}^{\mathbf{1}[y_{ij}=l]}
\]

\paragraph{M-step.} Update annotator error rates and class priors:
\[
\hat{\pi}_{jkl} = \frac{\sum_i P(z_i=k) \cdot \mathbf{1}[y_{ij}=l]}{\sum_i P(z_i=k)}
\qquad
\hat{p}_k = \frac{1}{N} \sum_i P(z_i = k)
\]

\begin{lstlisting}[style=python, caption={Dawid-Skene EM (simplified)}]
def dawid_skene_em(annotations, n_classes, n_iters=50):
    """
    annotations: dict {item_id: {annotator_id: label}}
    Returns: estimated true label posteriors T, annotator conf. matrices pi
    """
    N, K = len(annotations), n_classes
    items = list(annotations.keys())
    annotators = list(set(a for item in annotations.values() for a in item))
    M = len(annotators)

    # Initialize T with majority vote
    T = np.zeros((N, K))
    for i, item in enumerate(items):
        for v in annotations[item].values():
            T[i, v] += 1
        T[i] /= T[i].sum()

    pi = np.ones((M, K, K)) / K

    for _ in range(n_iters):
        # M-step: update pi
        for j, ann in enumerate(annotators):
            for k in range(K):
                for l in range(K):
                    num = sum(T[i,k] for i,item in enumerate(items)
                              if ann in annotations[item]
                              and annotations[item][ann] == l)
                    den = sum(T[i,k] for i,item in enumerate(items)
                              if ann in annotations[item])
                    pi[j,k,l] = num / (den + 1e-10)

        prior = T.mean(axis=0)

        # E-step: update T
        log_T = np.log(prior[np.newaxis,:] + 1e-10) * np.ones((N,K))
        for i, item in enumerate(items):
            for j, ann in enumerate(annotators):
                if ann in annotations[item]:
                    l = annotations[item][ann]
                    log_T[i] += np.log(pi[j,:,l] + 1e-10)

        log_T -= log_T.max(axis=1, keepdims=True)
        T = np.exp(log_T)
        T /= T.sum(axis=1, keepdims=True)

    return T, pi
\end{lstlisting}

% ============================================================================
\section{Framing Effects: Signal Detection Theory Framework}
% ============================================================================

Model annotation as a signal detection problem. The annotator observes item $x$ with latent ``harmfulness'' signal $s(x) \in \mathbb{R}$, and classifies by comparing to threshold $\theta$:

\[
\hat{y} = \begin{cases} 1 & s(x) > \theta \\ 0 & s(x) \leq \theta \end{cases}
\]

The annotation guideline framing shifts $\theta$ for individual annotators:
\[
\theta_A = \theta_0 + \Delta_A
\]

Inter-annotator agreement is a function of the variance in individual thresholds:
\[
\kappa \approx 1 - f(\sigma_\theta^2)
\]

Framing effects that reduce $\sigma_\theta^2$ (by more precisely specifying the threshold) directly increase $\kappa$. The ``reasonable person'' standard reduces idiosyncratic threshold variation at the cost of introducing shared cultural blind spots.

% ============================================================================
\section{Sample Size Planning for Annotation Projects}
% ============================================================================

\subsection{Statistical Power for Kappa Estimation}

To estimate kappa with margin of error $e$ at confidence level $\alpha$:

\[
N \geq \frac{z_{\alpha/2}^2 \cdot P_o(1-P_o)}{e^2 (1-P_e)^2}
\]

For typical annotation settings ($P_o \approx 0.80$, $P_e \approx 0.50$, $e = 0.05$, $\alpha = 0.05$):

\[
N \geq \frac{3.84 \times 0.16}{0.0025 \times 0.25} \approx 983 \text{ items}
\]

\subsection{Redundancy and Quality Trade-off}

Expected kappa improvement from $n_r$ annotators per item with majority-vote aggregation:

\[
\kappa(n_r) \approx \kappa_1 + (1 - \kappa_1) \cdot \left(1 - e^{-\lambda(n_r - 1)}\right)
\]

where $\kappa_1$ is single-annotator kappa and $\lambda \approx 0.3$. Increasing redundancy beyond $n_r = 5$ yields diminishing returns; the marginal quality gain from a 6th annotator is typically less than 0.01 kappa.

\begin{technote}[Empirical Constant $\lambda$]
The value $\lambda \approx 0.3$ is an empirical heuristic derived from annotation studies across multiple NLP and CV tasks; it is not derived from first principles. In practice, $\lambda$ varies by task difficulty and annotator pool. Treat the formula as a planning heuristic, not a precise prediction, and calibrate against pilot annotation data for any specific project.
\end{technote}

% ============================================================================
\section{Exercises and Assignments}
% ============================================================================

\begin{Exercise}[Cohen's Kappa Computation]
Two annotators classify 100 social media posts:
\begin{center}
\begin{tabular}{@{}lcc@{}}
\toprule
 & Ann.~2: T & Ann.~2: N \\
\midrule
Ann.~1: T & 20 & 15 \\
Ann.~1: N & 10 & 55 \\
\bottomrule
\end{tabular}
\end{center}
Compute $P_o$, $P_e$, $\kappa$, and a 95\% confidence interval. Interpret using the Landis-Koch scale and recommend whether to proceed with model training.
\end{Exercise}

\begin{Exercise}[Fleiss' Kappa Properties]
Three annotators each label 50 documents. Compute Fleiss' kappa and compare to average pairwise Cohen's kappa. Under what conditions do they diverge most significantly? Construct a 3-annotator example where average pairwise kappa is misleadingly optimistic.
\end{Exercise}

\begin{Exercise}[Dawid-Skene Simulation]
Simulate annotation data for 100 items with two expert annotators (error rate 5\%) and eight novice annotators (error rate 30\%). Compare accuracy of majority voting vs.\ Dawid-Skene EM vs.\ expert-only majority vote. How much does including novice annotators help or hurt under each method?
\end{Exercise}

\begin{Exercise}[Framing Effect Experiment]
Design an experiment to measure framing effects on a toxicity annotation task. Specify: two guideline versions, number of items, number of annotators per condition, statistical tests for kappa difference and positive label rate difference, and expected results based on the signal detection model.
\end{Exercise}

\begin{Exercise}[Annotation Budget Optimization]
Budget: \$5,000. Task: binary annotation, minimum 2,000 items. Expert: \$2.00/label, $\kappa_1 = 0.75$. Crowdworkers: \$0.05/label, $\kappa_1 = 0.50$. Design three annotation strategies. For each: estimate expected quality (using the redundancy formula), coverage, and total cost. Which strategy do you recommend, and under what conditions would you choose differently?
\end{Exercise}

\bigskip
\begin{center}
\rule{0.5\textwidth}{0.4pt}
\end{center}
\bigskip

\noindent\textit{This appendix provides the formal tools for quantifying annotation quality and diagnosing label problems described in Chapter~7. The kappa family of statistics, the Dawid-Skene model, and the signal detection framework together constitute the annotation quality toolkit for applied \hitl system development.}
